\documentclass[a4paper]{article}
% Second-round response; anonymized for double-blind review.
\def\papertitle{Correct-and-then-Detect: Adaptive Fault Detection through Data-driven Digital Twins and Transfer Learning}
\def\authors{Anonymous Authors}
\def\journal{ISA Transactions}
\def\paperid{ISATRANS-D-26-02059}
% Define title defaults if not defined by user
\providecommand{\lettertitle}{Response Letter of}
\providecommand{\papertitle}{Title}
\providecommand{\authors}{Authors}
\providecommand{\journal}{Journal}
\providecommand{\paperid}{--}

\usepackage[includeheadfoot,top=20mm, bottom=20mm, footskip=2.5cm]{geometry}

% Typography
\usepackage[T1]{fontenc}
\usepackage{times}
%\usepackage{mathptmx} % math also in times font
\usepackage{amssymb,amsmath,amsthm}
\usepackage{amsfonts}
\usepackage{array}
\usepackage{url}
\usepackage{microtype}
\usepackage[utf8]{inputenc}

% Misc
\usepackage{graphicx}
\graphicspath{{figs/}{../Figures/}} % Paths to the graph folder
\usepackage[hidelinks]{hyperref} %textopdfstring from pandoc
\usepackage{soul} % Highlight using \hl{}

% Include pdf pages
\usepackage{pdfpages}

% Table

\usepackage{booktabs}
\usepackage{adjustbox} % center large tables across textwidth by surrounding tabular with \begin{adjustbox}{center}
\renewcommand{\arraystretch}{1.5} % enlarge spacing between rows
\usepackage{caption}
\captionsetup[table]{skip=10pt} % enlarge spacing between caption and table
\usepackage{multirow}


% Section styles

\usepackage{titlesec}
\titleformat{\section}{\normalfont\large}{\makebox[0pt][r]{\bf \thesection.\hspace{4mm}}}{0em}{\bfseries}
\titleformat{\subsection}{\normalfont}{\makebox[0pt][r]{\bf \thesubsection.\hspace{4mm}}}{0em}{\bfseries}
\titlespacing{\subsection}{0em}{1em}{-0.3em} % left before after

% Paragraph styles

\setlength{\parskip}{0.6\baselineskip}%
\setlength{\parindent}{0pt}%

% Quotation styles

\usepackage{framed}
\let\oldquote=\quote
\let\endoldquote=\endquote
\renewenvironment{quote}{\begin{fquote}\advance\leftmargini -2.4em\begin{oldquote}   	\setlength{\parindent}{1em}}{\end{oldquote}\end{fquote}}

\usepackage{xcolor}
\newenvironment{fquote}
{\def\FrameCommand{
		\fboxsep=0.6em % box to text padding
		\fcolorbox{black}{white}}%
	% the "2" can be changed to make the box smaller
	\MakeFramed {\advance\hsize-2\width \FrameRestore}
	\begin{minipage}{\linewidth}
		}
		{\end{minipage}\endMakeFramed}

% Style for the revisions
\newcommand{\revisions}[1]{\textcolor{red}{\textit{#1}}}

% Number tables and figures as R1, R2, ... in this response letter
\renewcommand{\thetable}{R\arabic{table}}
\renewcommand{\thefigure}{R\arabic{figure}}

% Table styles

\let\oldtabular=\tabular
\let\endoldtabular=\endtabular
\renewenvironment{tabular}[1]{\begin{adjustbox}{center}\begin{oldtabular}{#1}}{\end{oldtabular}\end{adjustbox}}


% Shortcuts

%% Let textbf be both, bold and italic
%\DeclareTextFontCommand{\textbf}{\bfseries\em}

%% Add RC and AR to the left of a paragraph
%\def\RC{\makebox[0pt][r]{\bf RC:\hspace{4mm}}}
%\def\AR{\makebox[0pt][r]{AR:\hspace{4mm}}}

%% Define that \RC and \AR should start and format the whole paragraph 
\usepackage{suffix}
\long\def\RC#1\par{\makebox[0pt][r]{\bf RC:\hspace{4mm}}\textbf{#1}\par} %\RC
\WithSuffix\long\def\RC*#1\par{\textbf{{#1}}\par} %\RC*
\long\def\AR#1\par{\makebox[0pt][r]{AR:\hspace{10pt}}{#1}\par} %\AR
\WithSuffix\long\def\AR*#1\par{{#1}\par} %\AR*

%%%
%DIF PREAMBLE EXTENSION ADDED BY LATEXDIFF
%DIF UNDERLINE PREAMBLE %DIF PREAMBLE
\RequirePackage[normalem]{ulem} %DIF PREAMBLE
\RequirePackage{color} %DIF PREAMBLE
\providecommand{\DIFadd}[1]{{\protect\uwave{#1}}} %DIF PREAMBLE
\providecommand{\DIFdel}[1]{{\protect\sout{#1}}}                      %DIF PREAMBLE
%DIF SAFE PREAMBLE %DIF PREAMBLE
\providecommand{\DIFaddbegin}{} %DIF PREAMBLE
\providecommand{\DIFaddend}{} %DIF PREAMBLE
\providecommand{\DIFdelbegin}{} %DIF PREAMBLE
\providecommand{\DIFdelend}{} %DIF PREAMBLE
%DIF FLOATSAFE PREAMBLE %DIF PREAMBLE
\providecommand{\DIFaddFL}[1]{\DIFadd{#1}} %DIF PREAMBLE
\providecommand{\DIFdelFL}[1]{\DIFdel{#1}} %DIF PREAMBLE
\providecommand{\DIFaddbeginFL}{} %DIF PREAMBLE
\providecommand{\DIFaddendFL}{} %DIF PREAMBLE
\providecommand{\DIFdelbeginFL}{} %DIF PREAMBLE
\providecommand{\DIFdelendFL}{} %DIF PREAMBLE
%DIF END PREAMBLE EXTENSION ADDED BY LATEXDIFF

\begin{document}

% Make title
{\Large\bf \lettertitle}\\[1em]
{\huge \papertitle}\\[1em]
{\authors}\\
Manuscript submitted to {\it \journal}\\
Ms. Ref. No.:\paperid\\
\hrule

% Legend
\hfill {\bfseries RC:} \textbf{Reviewer Comment},\(\quad\) AR: \emph{Author Response}, \(\quad\square\) \emph{\revisions{Manuscript text}}

\section{Letter to the editor}
Dear editors,

Thank you for the continued assessment of our manuscript. We have addressed Reviewer \#3's remaining concern about the flow of the revised text. Both passages specifically identified by the reviewer have been rewritten, and the other revised passages have been reviewed to integrate explanations more naturally into the manuscript. The response below focuses on this remaining concern and quotes the relevant wording from the current revised manuscript. We have also made additional language, presentation, and consistency refinements elsewhere in the manuscript. Changes are highlighted in red in the anonymized manuscript.

With regards,

\authors

\newpage
\section*{Response to the Associate Editor}

\RC One of the reviewers still require further improvements before the paper could be accepted for publication. The authors should prepare a revised version of the paper addressing the remaining comments as thoroughly as possible.

\AR Thank you for your guidance. We have addressed Reviewer \#3's remaining concern and provide the specific revisions below. We have also made additional language and presentation refinements to improve the clarity and consistency of the manuscript.

\newpage
\section*{Response to the Editor in Chief}

\RC I concur with the recommendation from the handling editor. Reviewers report improvements; however, they are not unanimous that the latest version has addressed all the points raised by the reviewers. Should the authors opt to further revise and resubmit, they must emphasize the contributions of the paper beyond the current state-of-the-art and clearly establish its novelty. Additionally, they must provide a comprehensive and explicit point-by-point response to the reviewers in a separate document that observes the journal double-blind policy. In addition, the anonymized revised manuscript should be highlighted with all the changes made due to revision. It is important to note that revising the manuscript does not imply its ultimate acceptance. The revised version will be subjected to further scrutiny and review to ensure that it addresses all the issues raised and meets the standards expected for this publication.

\AR Thank you for the continued consideration of our manuscript. We have revised the contribution statement to more clearly present the novelty of the proposed framework and its relation to the limitations identified in existing data-driven digital twin methods. We have also integrated the methodological explanations more naturally into the surrounding discussion, as detailed in the response to Reviewer \#3. The response letter and revised manuscript remain anonymized, and the manuscript revisions are highlighted in red. Additional language, presentation, and consistency refinements have also been made.

\newpage
\section*{Response to Reviewer \#1}
\subsection*{Comment 1}

\RC The authors have made detailed revisions and improvements.  I have no more comments. It can now be accepted.

\AR Thank you for your careful review and positive assessment of the revised manuscript. We appreciate your recognition of the revisions and your recommendation for acceptance.

\newpage
\section*{Response to Reviewer \#2}
\subsection*{Comment 1}

\RC no comments.

\AR Thank you for reviewing the revised manuscript. We appreciate your time and consideration.

\newpage
\section*{Response to Reviewer \#3}
\subsection*{Comment 1}

\RC The authors' responses have largely addressed my concerns. However, one minor issue remains. Some revisions seem somewhat abrupt in the context. It is recommended to revise the relevant sections to make the writing flow more smoothly. For example:

\RC* Section 1: ``...It should be noted that the main contribution here is the integrated data-driven fault detection framework, including nominal behavior modeling, transfer-learning-based digital twin updating, and PROF-enhanced residual fault detection, rather than the performance of its individual building blocks like prediction accuracy of the digital twin. ...''

\RC* Section 3.3: ``...The freezing operation itself is not claimed as a new continual- or incremental-learning algorithm; it is the implementation used to maintain the nominal reference within the proposed detection framework. ...''

\RC* The examples listed above are not exhaustive. The authors are advised to carefully review all revised text.

\AR Thank you for this comment. We revised the two passages identified by the reviewer and reviewed the other revised sections for similar expressions. Several explanatory or repetitive statements were removed or shortened, and the corresponding passages were revised to read more naturally within the manuscript. The main revisions are summarized below.

\subsection*{1. Section 1: contribution statement}

\AR* We removed the sentence beginning ``It should be noted that the main contribution here is...'' and replaced it with the following statement:

\begin{quote}
\revisions{The proposed framework jointly addresses three coupled issues in data-driven digital-twin-based fault detection: representative nominal-data acquisition, adaptation of the nominal model under distribution shifts, and suppression of fault-induced contamination during rolling-horizon prediction. These functions are integrated into a unified workflow for maintaining reliable nominal references throughout fault detection.}
\end{quote}

\subsection*{2. Section 3.3: role of freezing and fine-tuning}

\AR* We removed the statement that ``the freezing operation itself is not claimed as a new continual- or incremental-learning algorithm'' and revised the corresponding text as follows:

\begin{quote}
\revisions{Within the proposed framework, layer-wise freezing and fine-tuning are used as the adaptation mechanism for updating the nominal model under verified healthy domain shifts.}
\end{quote}

\subsection*{3. Section 3.2.1: relation between task sampling and system responses}

\AR* We revised the description of how the sampled task configurations relate to the resulting system responses as follows:

\begin{quote}
\revisions{The Halton sequence is applied separately to the start and goal configuration spaces, and the sampled points are subsequently paired to form task configurations. These task configurations are converted into joint-space motions through inverse kinematics and trajectory planning, which in turn generate multivariate system responses. Accordingly, the proposed DoE provides representative coverage of the system responses generated by tasks distributed across the predefined task space.}
\end{quote}

\subsection*{4. Section 3.2.3: predictive-model interpretation}

\AR* We removed the previous explanatory statement about the interpretation of the model states and retained a direct definition of the model parameters and latent variables:

\begin{quote}
\revisions{Here, $\boldsymbol{\theta}_t$ denotes the trainable parameters of the deep predictive model, whereas its hidden variables are learned latent features rather than explicit physical states of the robot.}
\end{quote}

\AR* We also removed the redundant discussion of LSTM novelty. The implementation paragraph now states directly:

\begin{quote}
\revisions{In this work, $f_{\mathrm{DT}}(\cdot;\boldsymbol{\theta}_t)$ is implemented by a stacked long short-term memory (LSTM) network followed by a fully connected output layer. The network receives a $30\times18$ observation history and directly outputs the ten-step motor-temperature trajectory. LSTM gating equations are standard and are therefore omitted; the architecture and training settings required for reproduction are reported in Table~5. The proposed acquisition, adaptation, residual, and PROF procedures do not require an LSTM specifically.}
\end{quote}

\subsection*{5. Section 3.3: adaptation under healthy operating changes}

\AR* We revised the description of model updating to state directly when adaptation data are collected and how their health status is confirmed:

\begin{quote}
\revisions{When model updating is required, either for adaptation to a new device or to a new healthy operating stage of the current system, a batch of data is collected under controlled and confirmed healthy operating conditions and used for transfer-learning-based adaptation. This allows the updated digital twin to remain representative of the current nominal system behavior. In practical deployment, the health status of the adaptation data can be confirmed using maintenance information, additional monitoring signals, or controlled inspection.}
\end{quote}

\subsection*{6. Section 4.1: threshold calibration and transition to subsequent sections}

\AR* We shortened the description of the calibration-data requirement to:

\begin{quote}
\revisions{Model training and target-domain adaptation use healthy data only, while threshold calibration uses labeled fault-injected data only from the source-domain validation set.}
\end{quote}

\AR* The transition to Sections 4.2 and 4.3 was revised as follows:

\begin{quote}
\revisions{Beyond nominal prediction accuracy, reliable residual-based detection must also remain robust to two additional sources of mismatch: fault-induced contamination of historical inputs and cross-domain shifts in nominal behavior. These two issues are addressed in Sections 4.2 and 4.3, respectively.}
\end{quote}

\subsection*{7. Section 4.2: feedback-buffer behavior}

\AR* We removed the previous discussion of below-threshold abnormal observations and revised the description around the first detection time as follows:

\begin{quote}
\revisions{Before $t_d$, observations with residuals below the detection threshold remain unchanged in the feedback buffer. From $t_d$ onward, newly detected abnormal observations are replaced before the next rolling input is formed, thereby suppressing further recursive contamination.}
\end{quote}

\subsection*{8. Section 4.3: cross-domain threshold reuse}

\AR* We revised the description of cross-domain threshold reuse as follows:

\begin{quote}
\revisions{The source and target systems share the same monitored quantity, sensing resolution, sampling rate, preprocessing, residual definition, prediction horizon, and fault-decision rule. Under these consistent conditions, the source-calibrated threshold is reused in the target domain. The target nominal model is adapted using verified healthy data so that the target residual scale remains comparable with that of the source domain. These conditions hold for the two ArmPi robots considered in this study. Section 6.1 evaluates the domain shift and the alignment of target residuals after transfer learning, while Section 6.2.1 reports the threshold-sensitivity curve.}
\end{quote}

\subsection*{9. Section 5.4: scope of the thermal-fault validation}

\AR* We shortened the discussion of the experimental scope and retained the following statement:

\begin{quote}
\revisions{Real motor failures can additionally involve time-varying load, speed, ambient conditions, nonlinear heat transfer, controller intervention, and electrical, mechanical, or vibration signatures that are not represented by the first-order LPTN. Non-thermal, highly intermittent, or temperature-unobservable faults lie outside the current validation boundary. Accordingly, the present experiments validate the framework for progressive thermal anomalies, while validation using naturally occurring or safely reproduced physical faults remains necessary to establish its applicability to broader motor-failure scenarios.}
\end{quote}

\AR* In addition, we made minor language, terminology, figure, table, and cross-reference corrections throughout the manuscript for consistency.

\newpage
\section*{Response to Reviewer \#4}
\subsection*{Comment 1}

\RC Accept

\AR Thank you for your review and recommendation for acceptance. We appreciate your time and consideration.

\end{document}
